\subsection{An analytic completion of the all-Picard certificate}
\label{subsec:analytic-all-picard}

The local example in \cref{prop:local-cusp-picard-obstruction} shows that
the cusp neighbourhood alone cannot annihilate normalized central
restrictions.  We now close the stronger analytic route by retaining the
order-three and order-four fillings.  The proof has two independent parts.
First, logarithms of the actual \(\mathcal O^*\)-factors give the boundary
cocycle and its two finite-filling norm identities.  Second, a
candidate-specific support-function calculation reads every possible
central restriction directly from the mixed Appell--Humbert block on a
nearby torus.  No global Leray differential, nearby-cycle specialization,
or Chern-class evaluation on a cusp bounding chain occurs.

\begin{proposition}[analytic Appell--Humbert boundary cocycle]
\label{prop:analytic-ah-boundary}
\Verified\ Let \(P\subset B^\circ\) be the compact pair-of-pants core,
orient its boundary meridians clockwise, and retain the zero-section lifts
of \cref{prop:geometric-pi1}.  Thus
\[
 \rho_1\rho_2\rho_0=1,
 \qquad
 \widehat\rho_1\widehat\rho_2\widetilde a_0=1,
\tag{5.17-A1}
\]
and the second relation has no terminal vertical translation.  For
\(L\in\operatorname{Pic}(X)\), let
\[
 E=c_1(L|_F)\in\operatorname{Alt}^2(\Lambda,\mathbb Z)
   =\bigwedge^2V
\]
be its fibrewise Appell--Humbert form.  Put
\(V_{\mathbb R}=\operatorname{Hom}(\Lambda,\mathbb R)\), and use
\[
 A_g=(T_g^{-1})^{\mathsf t},
 \qquad
 \langle T_g\alpha,\lambda\rangle
 =\langle\alpha,A_g^{-1}\lambda\rangle.
\tag{5.17-A2}
\]
There is a crossed homomorphism, unique up to a principal crossed
homomorphism,
\[
 d:\pi_1(P)\longrightarrow V_{\mathbb R},
 \qquad
 d(gh)=d(g)+T_gd(h),
 \qquad
 d(g)\sim d(g)+(T_g-I)a.
\tag{5.17-A3}
\]
If \(d_i=d(\rho_i)\), then
\[
 d_1+T_1d_2+T_1T_2d_0=0,
 \qquad
 \bar d_1+\bar d_2+\bar d_0=0
 \quad\text{in }V_G\otimes\mathbb R.
\tag{5.17-A4}
\]
For the finite filling of order \(m_j\), \(j=1,2\),
\[
 N_{T_j}d_j=\iota_{v_j}E,
 \qquad
 m_j\bar d_j=[\iota_{v_j}E],
 \qquad
 N_{T_j}=I+T_j+\cdots+T_j^{m_j-1},
\tag{5.17-A5}
\]
where \((\iota_vE)(\lambda)=E(v,\lambda)\).  At the cusp,
\[
 \langle d_0,\widehat\delta\rangle=0,
 \qquad
 \bar d_0=0\quad\text{in }V_G\otimes\mathbb R.
\tag{5.17-A6}
\]
These identities are obtained from holomorphic transition factors and the
argument principle.
\end{proposition}

\begin{proof}
Choose an open Riemann-surface neighbourhood \(P^{\rm an}\Subset B^\circ\)
which deformation retracts onto \(P\), and pull the period family to its
universal cover.  Before quotienting by the fibre lattice, the resulting
space is
\[
 \widetilde{\mathcal J}
 =\widetilde P^{\rm an}\times\mathbb C^2.
\]
It is Stein and contractible, so the pullback of \(L\) is holomorphically
trivial.  In a nowhere-zero frame, write the factor belonging to a deck
transformation \(\alpha\) as
\[
 a_\alpha\in\mathcal O^*(\widetilde{\mathcal J}),
 \qquad
 a_{\alpha\beta}(x)=a_\alpha(\beta x)a_\beta(x).
\tag{5.17-A7}
\]
The space is simply connected, hence \(a_\alpha=e^{2\pi i h_\alpha}\) for
a holomorphic logarithm.  Therefore
\[
 c(\alpha,\beta)
 =h_\alpha(\beta x)+h_\beta(x)-h_{\alpha\beta}(x)\in\mathbb Z
\tag{5.17-A8}
\]
is independent of \(x\), and associativity makes \(c\) an integral group
two-cocycle.

Let \(\tau_\lambda\) denote fibre translation by
\(\Pi(z)\lambda\), \(\lambda\in\Lambda\), and let \(\widehat g\) be the
zero-section lift of a base deck transformation.  Period equivariance gives
\[
 \widehat g\,\tau_\lambda\,\widehat g^{-1}
 =\tau_{A_g\lambda}.
\tag{5.17-A9}
\]
On fibre translations, the commutator of (5.17-A8) is
\[
 E(\lambda,\mu)
 =c(\tau_\lambda,\tau_\mu)-c(\tau_\mu,\tau_\lambda).
\tag{5.17-A10}
\]
The four oriented sides of the \((\lambda,\mu)\)-parallelogram identify
this integer with the exponential connecting class of \(L|_F\); this is
the Appell--Humbert convention used in
(5.17av)--(5.17ax), and agrees with the standard factor-of-automorphy
calculation \cite[Chapter~1, \S2, pp.~5--8, especially
Theorem~2.1, p.~7]{BirkenhakeLange1999}%
\cite[equations~(3.3)--(3.10), pp.~3--7]{SchnellLecture3}.
Conjugating the commutator by \(\widehat g\) yields
\[
 E(A_g\lambda,A_g\mu)=E(\lambda,\mu).
\tag{5.17-A11}
\]

Push the logarithmic central extension defined by \(c\) out along
\(\mathbb Z\hookrightarrow\mathbb R\), and denote its central element of
weight \(r\) by \(\mathbf c(r)\).  Choose lifts \(U_\lambda\) of the fibre
translations satisfying the real Appell--Humbert normalization
\[
 U_\lambda U_\mu
 =\mathbf c\!\left(\frac12E(\lambda,\mu)\right)U_{\lambda+\mu}.
\tag{5.17-A12}
\]
To see that this normalization exists, start with arbitrary lifts.  The
difference between their real cocycle and \(E/2\) has zero commutator, so
it defines an abelian extension of the free abelian group \(\Lambda\) by
\(\mathbb R\).  Choose lifts of an integral basis and extend
multiplicatively; this splits that abelian extension and gives (5.17-A12).
If \(U'_\lambda\) is a second normalized system, then
\(U'_\lambda U_\lambda^{-1}\) is central, say
\(\mathbf c(\langle a,\lambda\rangle)\).  Comparing the two copies of
(5.17-A12) proves additivity in \(\lambda\), so
\(a\in V_{\mathbb R}\).  Thus every two normalized systems differ by exactly
the character change used below.

For \(g\in\pi_1(P)\), choose a lift \(G_g\) of \(\widehat g\) and define
\(d(g)\in V_{\mathbb R}\) by
\[
 G_gU_{A_g^{-1}\lambda}G_g^{-1}
 =\mathbf c\bigl(\langle d(g),\lambda\rangle\bigr)U_\lambda.
\tag{5.17-A13}
\]
The exponent is additive in \(\lambda\): conjugate (5.17-A12) and use
(5.17-A11).  The lifts of \(gh\) and \(G_gG_h\) differ centrally, so
\[
\begin{aligned}
 \langle d(gh),\lambda\rangle
 &=\langle d(g),\lambda\rangle
   +\langle d(h),A_g^{-1}\lambda\rangle\\
 &=\langle d(g)+T_gd(h),\lambda\rangle.
\end{aligned}
\]
This proves the crossed-homomorphism law.  Replacing
\(U_\lambda\) by
\(\mathbf c(\langle a,\lambda\rangle)U_\lambda\) changes (5.17-A13) by
\((T_g-I)a\).  Replacing \(G_g\) by a central multiple does not change
(5.17-A13), so this proves the stated independence and uniqueness.  Applying
the crossed law to
(5.17-A1) proves the first identity in (5.17-A4); passage to coinvariants
proves the second.

For the order-\(m_j\) filling, distinguish every covering level.  Put
\[
 \widetilde J_j=\Delta_j\times\mathbb C^2,\qquad
 r_j:\widetilde J_j\longrightarrow J'_j
\]
for the lattice quotient, whose deck translations are \(\tau_\lambda\).
On the vector cover define
\[
 Q_{j,\log}:=\widetilde g_j^{\,v_j/m_j},\qquad
 Q_{j,\mathrm{lin}}:=\widetilde g_j^{\,0}.
\]
They satisfy
\[
\begin{aligned}
 Q_{j,\log}\tau_\lambda Q_{j,\log}^{-1}
  &=\tau_{A_j\lambda},&
 (Q_{j,\log})^{m_j}&=\tau_{v_j},\\
 Q_{j,\mathrm{lin}}\tau_\lambda Q_{j,\mathrm{lin}}^{-1}
  &=\tau_{A_j\lambda},&
 (Q_{j,\mathrm{lin}})^{m_j}&=1.
\end{aligned}
\tag{5.17-A14}
\]
The positive sign is fixed by
\(s_j\circ g_j=e^{-2\pi i/m_j}s_j\) and (2.8f): after exactly
\(m_j\) clockwise iterations, the logarithmic vector-cover lift translates
by \(+v_j\).

Let \(q_{j,\log},q_{j,\mathrm{lin}}\) be the descended automorphisms of \(J'_j\),
and put \(J_j^{\prime *}=J'_j|_{\Delta_j^*}\).  Both descended automorphisms have
order dividing \(m_j\), because \(\tau_{v_j}\) is trivial after the lattice
quotient; each covers \(g_j\), whose order is exactly \(m_j\), so each
descended order is exactly \(m_j\).
The branch-independent torus translation
\[
 h_j:J_j^{\prime *}\longrightarrow J_j^{\prime *}
\]
constructed in (2.9d) satisfies, on this quotient level,
\[
 h_jq_{j,\log}h_j^{-1}=q_{j,\mathrm{lin}}.
\]
Define
\[
\begin{aligned}
 p_{j,\log}:J_j^{\prime *}&\longrightarrow
  N_j^*:=J_j^{\prime *}/\langle q_{j,\log}\rangle,\\
 p_{j,\mathrm{lin}}:J_j^{\prime *}&\longrightarrow
 J_j^*:=J_j^{\prime *}/\langle q_{j,\mathrm{lin}}\rangle .
\end{aligned}
\]
The conjugacy induces a biholomorphism
\[
 \bar h_j:N_j^*\longrightarrow J_j^*,\qquad
 p_{j,\mathrm{lin}}\circ h_j=\bar h_j\circ p_{j,\log}.
\]
The gluing of \(X\) identifies
\(L|_{N_j^*}\simeq\bar h_j^*(L|_{J_j^*})\); hence the exactly typed
line-bundle square gives
\[
 p_{j,\log}^*(L|_{N_j^*})
 \simeq h_j^*p_{j,\mathrm{lin}}^*(L|_{J_j^*}).
\tag{5.17-A14a}
\]

For the factor extension return to the filled vector cover.  Let
\[
 P_{j,\log}:J'_j\longrightarrow N_j,\qquad
 \pi_j^{\mathrm{vec}}:=P_{j,\log}\circ r_j:
 \widetilde J_j\longrightarrow N_j .
\]
Here \(P_{j,\log}\) is the full quotient and
\(P_{j,\log}|_{J_j^{\prime *}}=p_{j,\log}\).
The deck group is
\[
 \Gamma_j^{\log}
 =\langle\tau_\lambda,Q_{j,\log}:\lambda\in\Lambda\rangle.
\]
The pullback \((\pi_j^{\mathrm{vec}})^*(L|_{N_j})\) is trivial on the
contractible Stein manifold \(\widetilde J_j\); its factors define an
integral central extension, which we push out along
\(\mathbb Z\hookrightarrow\mathbb R\).  Over \(\Delta_j^*\), the
line-bundle square (5.17-A14a) identifies each fibre with the global factor
extension by a translation, so its commutator form is \(E\).  That form is
integer-valued and locally constant on the connected disc \(\Delta_j\);
therefore it is \(E\) on the filled cover as well.  Choose vertical lifts
\(U_\lambda^{(j)}\) normalized by
\[
 U_\lambda^{(j)}U_\mu^{(j)}
 =\mathbf c\!\left(\frac12E(\lambda,\mu)\right)
  U_{\lambda+\mu}^{(j)}.
\]
For \(\alpha\in\Gamma_j^{\log}\), define \(A_\alpha,T_\alpha\) by
\[
 \alpha\tau_\lambda\alpha^{-1}=\tau_{A_\alpha\lambda},\qquad
 T_\alpha=(A_\alpha^{-1})^{\mathsf t}.
\]
For a lift \(\widetilde G_\alpha\), define
\(\widetilde d_j(\alpha)\in V_{\mathbb R}\) through
\[
 \widetilde G_\alpha U_{A_\alpha^{-1}\lambda}^{(j)}
 \widetilde G_\alpha^{-1}
 =\mathbf c\!\left(\langle\widetilde d_j(\alpha),\lambda\rangle\right)
  U_\lambda^{(j)}.
\tag{5.17-A14c}
\]
Indeed, before naming \(\widetilde d_j(\alpha)\), the quotient of the two
sides of (5.17-A14c) is central because both cover the same translation.
Conjugating the normalized product and using
\[
 E(A_\alpha\lambda,A_\alpha\mu)=E(\lambda,\mu)
\]
shows that its exponent \(r_\alpha(\lambda)\) satisfies
\[
 r_\alpha(\lambda+\mu)=r_\alpha(\lambda)+r_\alpha(\mu).
\]
Thus it is the pairing with a unique element of \(V_{\mathbb R}\), as
asserted.  Since \(A_{\alpha\beta}=A_\alpha A_\beta\), choose the lift of
\(\alpha\beta\) to be \(\widetilde G_\alpha\widetilde G_\beta\); any other
choice differs centrally.  Then
\[
\begin{aligned}
 &\widetilde G_\alpha\widetilde G_\beta
 U_{A_\beta^{-1}A_\alpha^{-1}\lambda}^{(j)}
 \widetilde G_\beta^{-1}\widetilde G_\alpha^{-1}\\
 &\quad=
 \mathbf c\!\left(
  \langle\widetilde d_j(\beta),A_\alpha^{-1}\lambda\rangle
  +\langle\widetilde d_j(\alpha),\lambda\rangle\right)
 U_\lambda^{(j)}.
\end{aligned}
\]
The defining dual action
\(\langle T_\alpha v,\lambda\rangle
=\langle v,A_\alpha^{-1}\lambda\rangle\) therefore proves the crossed law
\[
 \widetilde d_j(\alpha\beta)
 =\widetilde d_j(\alpha)+T_\alpha\widetilde d_j(\beta).
\tag{5.17-A14d}
\]

Put
\[
 \widetilde J_j^*:=\Delta_j^*\times\mathbb C^2,
 \qquad
 r_j^*:=r_j|_{\widetilde J_j^*}:
 \widetilde J_j^*\longrightarrow J_j^{\prime *}
\]
and define
\[
\begin{aligned}
 \Phi_{j,\log}
  &:=p_{j,\log}\circ r_j^*,\\
 \Phi_{j,\mathrm{glob}}
  &:=p_{j,\mathrm{lin}}\circ h_j\circ r_j^*
   =\bar h_j\circ\Phi_{j,\log}.
\end{aligned}
\]
Both maps are invariant under \(\Gamma_j^{\log}\).  This is immediate for
every \(\tau_\lambda\), while
\[
\begin{aligned}
 \Phi_{j,\log}Q_{j,\log}
 &=p_{j,\log}q_{j,\log}r_j^*
  =\Phi_{j,\log},\\
 \Phi_{j,\mathrm{glob}}Q_{j,\log}
 &=p_{j,\mathrm{lin}}h_jq_{j,\log}r_j^*
  =p_{j,\mathrm{lin}}q_{j,\mathrm{lin}}h_jr_j^*
  =\Phi_{j,\mathrm{glob}}.
\end{aligned}
\]
Thus \(\widetilde J_j^*\) is the same \(\Gamma_j^{\log}\)-cover of the two
biholomorphic quotients.  Pulling the gluing isomorphism to this cover gives
\[
 \Theta_j:
 \Phi_{j,\log}^*(L|_{N_j^*})
 \xrightarrow{\ \sim\ }
 \Phi_{j,\mathrm{glob}}^*(L|_{J_j^*}).
\]
The filled factor extension restricts to this cover, and the global factor
extension pulls back to it through the fixed boundary whisker.  Because
\(\Theta_j\) intertwines descent, it induces an isomorphism of the resulting
pushed-out \(\Gamma_j^{\log}\)-factor extensions which sends
\(\mathbf c(r)\) to \(\mathbf c(r)\) for every \(r\in\mathbb R\).

Let \(\gamma_j^J:[0,1]\to J_j^*\) be the zero-section image of the
unwhiskered clockwise boundary meridian, and let \(\omega_j\) be the fixed
zero-section whisker from the global basepoint to \(\gamma_j^J(0)\).  Thus
the based total-space loop retained in (5.17-A1) is
\[
 \widehat\rho_j=\omega_j*\gamma_j^J*\omega_j^{-1}.
\]
Choose \(x_j\in J_j^{\prime *}\) such that
\(p_{j,\mathrm{lin}}(h_j(x_j))=\gamma_j^J(0)\), put
\(x_{j,\mathrm{lin}}=h_j(x_j)\), and let
\(\widetilde\gamma_{j,\mathrm{lin}}\) be the
\(p_{j,\mathrm{lin}}\)-lift of \(\gamma_j^J\) beginning at
\(x_{j,\mathrm{lin}}\).  For the chosen clockwise generator,
\[
 \widetilde\gamma_{j,\mathrm{lin}}(1)
 =q_{j,\mathrm{lin}}(x_{j,\mathrm{lin}}).
\]
Set
\[
 \gamma_j^N:=\bar h_j^{-1}\circ\gamma_j^J,
 \qquad
 \widetilde\gamma_{j,\log}
 :=h_j^{-1}\circ\widetilde\gamma_{j,\mathrm{lin}}.
\]
Then the quotient square gives the point-valued identities
\[
\begin{aligned}
 p_{j,\log}\circ\widetilde\gamma_{j,\log}
  &=\gamma_j^N,\\
 \widetilde\gamma_{j,\log}(0)&=x_j,\\
 \widetilde\gamma_{j,\log}(1)&=q_{j,\log}(x_j).
\end{aligned}
\]
Choose \(\widetilde x_j\in\widetilde J_j^*\) with
\(r_j^*(\widetilde x_j)=x_j\), and let
\(\widetilde\gamma_{j,\mathrm{vec}}\) be the \(r_j^*\)-lift of
\(\widetilde\gamma_{j,\log}\) beginning at \(\widetilde x_j\).  On
adjacent sheets choose compatible logarithm branches \(\ell,\ell'\) with
\[
 \ell'(w)=\ell(g_j^{-1}w)-\frac{2\pi i}{m_j}.
\]
The calculation in (2.9e) gives the point-valued endpoint formula
\[
 \widetilde\gamma_{j,\mathrm{vec}}(1)
 =Q_{j,\log}(\widetilde x_j),
\]
excluding both \(Q_{j,\log}^{-1}(\widetilde x_j)\) and every
\(\tau_\lambda Q_{j,\log}(\widetilde x_j)\) with \(\lambda\ne0\).
Fibre loops remain \(\tau_\lambda\), since \(h_j\) is a fibre translation.

The factor of \(L\) along the fixed lift of \(\omega_j\) identifies the
target factor over the local boundary segment with the globally based
factor.  Let \(G_{\rho_j}^{(j)}\) be the target lift of the local boundary
deck element whose whiskered image is the chosen global lift
\(G_{\rho_j}\).  Thus it covers \(\widehat\rho_j\) after whiskering and,
by (5.17-A13), satisfies
\[
 G_{\rho_j}^{(j)}U_{A_j^{-1}\lambda}(G_{\rho_j}^{(j)})^{-1}
 =\mathbf c(\langle d_j,\lambda\rangle)U_\lambda,
 \qquad d_j=d(\rho_j).
\]
Consequently, for some
\(a_j\in V_{\mathbb R}\) and \(b_j\in\mathbb R\),
\[
 \Theta_{j,*}(U_\lambda^{(j)})
 =\mathbf c(\langle a_j,\lambda\rangle)U_\lambda,
 \qquad
 \Theta_{j,*}(\widetilde G_{Q_{j,\log}})
 =\mathbf c(b_j)G_{\rho_j}^{(j)}.
\]
Here \(U_\lambda\) is the normalized global system in (5.17-A12), and
\(\mathbf c(b_j)\) disappears under conjugation.
Therefore
\[
\begin{aligned}
 \langle\widetilde d_j(Q_{j,\log}),\lambda\rangle
 +\langle a_j,\lambda\rangle
 &=\langle d_j,\lambda\rangle
   +\langle a_j,A_j^{-1}\lambda\rangle\\
 &=\langle d_j+T_ja_j,\lambda\rangle,
\end{aligned}
\]
so
\[
 \widetilde d_j(Q_{j,\log})=d_j+(T_j-I)a_j.
\tag{5.17-A14e}
\]

For a translation choose
\(\widetilde G_{\tau_v}=U_v^{(j)}\).  The filled normalization above gives
\[
 U_v^{(j)}U_\lambda^{(j)}(U_v^{(j)})^{-1}
 =\mathbf c(E(v,\lambda))U_\lambda^{(j)},
\]
and hence
\[
 \widetilde d_j(\tau_v)=\iota_vE,\qquad
 (\iota_vE)(\lambda)=E(v,\lambda).
\tag{5.17-A14f}
\]
Here
\[
 A_{Q_{j,\log}}=A_j,\qquad T_{Q_{j,\log}}=T_j.
\]
Since \(N_{T_j}(T_j-I)=T_j^{m_j}-I=0\), equations
(5.17-A14), (5.17-A14d), and (5.17-A14e) give
\[
\begin{aligned}
 N_{T_j}d_j
 &=N_{T_j}\widetilde d_j(Q_{j,\log})\\
 &=\widetilde d_j((Q_{j,\log})^{m_j})\\
 &=\widetilde d_j(\tau_{v_j})\\
 &=\iota_{v_j}E.
\end{aligned}
\tag{5.17-A15}
\]
Changing a line-bundle frame gives an isomorphic factor extension; changing
\(\widetilde G_\alpha\) by a central factor does not change conjugation; and
every second normalized vertical system differs by a character, changing
\(\widetilde d_j(\alpha)\) by \((T_\alpha-I)a\).  Finally,
\(\ell\mapsto\ell+2\pi in\) changes a vector-cover lift of \(h_j\) by
\[
 H_j\longmapsto H'_j=\tau_{nv_j}H_j.
\]
At the factor-extension level this means
\[
 \Theta'_{j,*}
 =\operatorname{Ad}(U_{nv_j})\circ\Theta_{j,*}.
\]
Consequently the comparison character changes by
\(n\iota_{v_j}E\), because
\[
 U_{nv_j}
 \bigl(\mathbf c(\langle a_j,\lambda\rangle)U_\lambda\bigr)
 U_{nv_j}^{-1}
 =
 \mathbf c(\langle a_j+n\iota_{v_j}E,\lambda\rangle)U_\lambda.
\]
Since \(A_jv_j=v_j\) and \(E\) is \(A_j\)-invariant,
\[
 \langle T_j(\iota_{v_j}E),\lambda\rangle
 =E(v_j,A_j^{-1}\lambda)
 =E(A_jv_j,\lambda)
 =E(v_j,\lambda),
\]
so this branch change leaves both (5.17-A14e) and (5.17-A15) unchanged.
The translation \(\tau_{nv_j}\) also commutes with
\(Q_{j,\mathrm{lin}}\) on the vector cover because \(A_jv_j=v_j\); on the
lattice quotient it is already the identity.  Hence it changes only the
central lift choice there.
The descended \(h_j,\bar h_j\), and the quotient square are branch
independent.  Thus (5.17-A15) is choice-independent and proves (5.17-A5),
including its sign.

It remains to obtain (5.17-A6) analytically.  Since
\(A_0\widehat\delta=\widehat\delta\), equation (5.17-A13) says that
\[
 [G_0,U_{\widehat\delta}]
 =\mathbf c\bigl(\langle d_0,\widehat\delta\rangle\bigr).
\tag{5.17-A16}
\]
This logarithmic commutator is the clutching winding with the base loop
first.  Thus, if \(b\) and \(f\) denote the oriented base and fibre loops,
respectively,
\[
 \operatorname{wind}_{b\wedge f}(L)
  =\langle d_0,\widehat\delta\rangle,
 \qquad
 \operatorname{wind}_{\xi_{\widehat\delta}}(L)
 =-\langle d_0,\widehat\delta\rangle.
\tag{5.17-A17}
\]
The second equality uses the fibre-first convention
\(\xi_{\widehat\delta}=[S^1_{\widehat\delta}]\times[\rho_0]\), hence the
single sign change.
Use the actual unimodular cusp chart
\[
 (u,x_1,x_2)=\left(\frac{t}{x_1x_2},x_1,x_2\right),
 \qquad t=ux_1x_2.
\tag{5.17-A18}
\]
For \(0<\rho\ll1\), the boundary parametrization
\[
 (\theta,\phi)\longmapsto
 \bigl(\rho e^{i(\theta-\phi)},1,e^{i\phi}\bigr),
 \qquad t=\rho e^{i\theta},
\tag{5.17-A19}
\]
has base coordinate \(\theta\) and fibre translation
\(\widehat\delta\) in the \(\phi\)-direction.  The latter fixed-base loop
does \emph{not} itself bound the \(x_2\)-disc.  To identify the correct
meridian, put \(\alpha=\theta-\phi\), let \(a_u\) be the positively
oriented \(u\)-longitude, and let \(m_{x_2}\) be the positively oriented
\(x_2\)-meridian at fixed \(u\).  In the basis
\((a_u,m_{x_2})\), the base and fibre loops are
\[
 b=a_u,
 \qquad
 f=m_{x_2}-a_u,
 \qquad
 b\wedge f=a_u\wedge m_{x_2}.
\tag{5.17-A19a}
\]
The change-of-basis determinant is \(+1\), so this preserves the
base-first orientation in (5.17-A17).  Now \(m_{x_2}\) does bound the
analytic disc
\[
 z\longmapsto(\rho e^{i\alpha_0},1,z),
 \qquad |z|\le1.
\tag{5.17-A19b}
\]

Choose a thin annulus \(A_u\) about \(|u|=\rho\), a disc \(D_{x_1}\)
about \(1\), and a disc \(D_{x_2}\) containing \(\{|x_2|\le1\}\), small
enough that
\[
 A_u\times D_{x_1}\times D_{x_2}
 \subset\{|ux_1x_2|<\epsilon\}.
\tag{5.17-A19c}
\]
This neighbourhood is Stein and retracts onto \(S^1\), so its second
integral cohomology vanishes.  Cartan's theorem~B and the exponential
sequence therefore give trivial Picard group
\cite[\S2, fundamental theorem, p.~81]{Cartan1957}%
\cite[Proposition~6.28 and equations~(6.35)--(6.36),
pp.~180--182]{Lee2024}.  Pull \(L|_{N_0}\) back to
the cusp cover and choose one holomorphic frame on this entire
neighbourhood.  Compute the clutching cocycle with this extending frame on
both sides of a cut in the \(u\)-longitude.  Its transition unit on
\(D_{x_2}\) is then identically \(F=1\), so by (5.17-A19a) its
base--fibre commutator winding is
\[
 \frac1{2\pi i}\int_{|x_2|=1}\frac{dF}{F}=0.
\tag{5.17-A20}
\]
This computes the same alternating logarithmic cocycle as (5.17-A16).
Indeed, changing from the global Appell--Humbert frame to the extending
chart frame changes \(c\) by a group coboundary.  Put
\[
 g=\widehat g_0,
 \qquad
 h=\tau_{\widehat\delta}
\]
in the underlying deck group.  Equation (5.17-A9) and
\(A_0\widehat\delta=\widehat\delta\) give \(gh=hg\).  If
\(c'=c+\delta\beta\) is the cocycle after a frame change, then
\[
 \bigl[c'(g,h)-c'(h,g)\bigr]
 -\bigl[c(g,h)-c(h,g)\bigr]
 =\beta(hg)-\beta(gh)=0.
\]
Thus the alternating cocycle is frame-invariant without assuming that its
central-extension lifts commute.  Equations (5.17-A17), (5.17-A19a), and
(5.17-A20) therefore give
\(\langle d_0,\widehat\delta\rangle=0\).  This is a calculation in
\(\mathcal O^*\), not an evaluation of a Chern class on a bounding chain.
Finally, the matrices \(A_1,A_2\) fix
\(\widehat\delta\), so the evaluation is unchanged by every principal
change in (5.17-A3).  Since
\(V_G=\mathbb Z[\delta]\) and
\(\delta(\widehat\delta)=1\), its vanishing is equivalent to
\(\bar d_0=0\).  This proves (5.17-A6).
\end{proof}

\begin{lemma}[the toric support function detects central restriction]
\label{lem:analytic-support-specialization}
\Verified\ Let \(M\in\operatorname{Pic}(N_0)\), let
\(\pi:Y_\epsilon\to N_0\) be the cusp covering, and let
\(\nu:E_0\to W\) be the normalization.  If the Appell--Humbert form of
\(M|_{F_t}\), \(t\ne0\), is \(E_M\), then the three independent degrees
of \(\nu^*(M|_W)\) determine an integral symmetric form
\(B:\mathbb Z^2\times\mathbb Z^2\to\mathbb Z\) satisfying
\[
 E_M\bigl((\lambda,0),(0,m)\bigr)
 =-B(B_0\lambda,m)
 \qquad(\lambda,m\in\mathbb Z^2).
\tag{5.17-A21}
\]
Consequently
\[
 E_M=0
 \quad\Longrightarrow\quad
 \nu^*(M|_W)\simeq\mathcal O_{E_0}.
\tag{5.17-A22}
\]
\end{lemma}

\begin{proof}
Put \(\mathcal M=\pi^*M\).  The identity
\(\pi\circ\Psi_\lambda=\pi\) gives \(\mathcal M\) its genuine
\(\Psi\)-equivariant descent structure.  In the notation of
\cref{prop:local-cusp-picard-obstruction}, order the six rays in the star
of the origin as
\[
 (r_1,r_2,r_3,r_4,r_5,r_6)
 =(e_1,e_2,e_2-e_1,-e_1,-e_2,e_1-e_2),
 \qquad C_i=E_0\cap E_{r_i}.
\tag{5.17-A23}
\]
The quotient identifies opposite conductor branches.  Therefore, for
\(A=\mathcal M|_{E_0}=\nu^*(M|_W)\),
\[
 d_i=\deg_{C_i}A,
 \qquad
 (\deg_{C_1}A,\ldots,\deg_{C_6}A)
 =(d_1,d_2,d_3,d_1,d_2,d_3).
\tag{5.17-A24}
\]

Define the integer-valued quadratic function
\[
 a(x,y)
 =(d_2+d_3)\binom{x}{2}+d_3xy+(d_1+d_3)\binom{y}{2}.
\tag{5.17-A25}
\]
Its polarization is
\[
 B(v,w)=a(v+w)-a(v)-a(w),
 \qquad
 [B]=
 \begin{pmatrix}
  d_2+d_3&d_3\\ d_3&d_1+d_3
 \end{pmatrix}.
\tag{5.17-A26}
\]
The locally finite divisor
\[
 D_a=\sum_{v\in\mathbb Z^2}a(v)E_v
\tag{5.17-A27}
\]
is Cartier: on every maximal chart, the three height-one ray generators
form an integral basis, so the three assigned coefficients are the values
of one integral character.  In the standard toric convention the associated
support function has ray values \(-a(v)\); the divisor coefficients here are
the values \(a(v)\), consistent with
\(\operatorname{div}(\chi^m)=\sum_\rho\langle m,u_\rho\rangle D_\rho\)
\cite[\S3.3, pp.~60--63; \S3.4, pp.~63--72]{Fulton1993}%
\cite[Propositions~4.1.1--4.1.2, p.~171; Theorems~4.2.8,
p.~179 and~4.2.12, p.~182]{CoxLittleSchenck2011}.
This is the same chartwise Cartier argument as
in (5.17ab), and it applies to an arbitrary integer-valued \(a\).

Direct substitution gives
\[
 (a(r_1),\ldots,a(r_6))
 =(0,0,d_2,d_2+d_3,d_1+d_3,d_1).
\tag{5.17-A28}
\]
Using \(C_i^2=-1\) and \(C_iC_{i\pm1}=1\), the degree vector of
\(\mathcal O(D_a)|_{E_0}\) is exactly the vector in (5.17-A24).  More
generally, restrict the principal identity
\(\operatorname{div}(t)=\sum_wE_w\) to \(E_v\).  Since \(E_v\) meets
precisely \(E_{v+r_i}\), it gives
\[
 \mathcal O(D_a)|_{E_v}
 \simeq
 \mathcal O_{E_v}\!\left(
   \sum_{i=1}^6(a(v+r_i)-a(v))C_{v,i}
 \right).
\tag{5.17-A29}
\]
The coefficient difference is
\[
 a(v+r_i)-a(v)=a(r_i)+B(v,r_i).
\tag{5.17-A30}
\]
The second summand in (5.17-A30) is the divisor of the toric character
with exponent \(B(v,-)\).  Thus every restriction in (5.17-A29) has the
same six degrees as \(\mathcal O(D_a)|_{E_0}\).

Because \(B_0\bar\Lambda=\mathbb Z^2\), the deck action is transitive on
the components \(E_v\).  The genuine descent structure on \(\mathcal M\)
therefore transports (5.17-A24), with its ordered six boundary directions,
to every \(E_v\).  It follows that
\[
 \mathcal R:=\mathcal M\otimes\mathcal O(-D_a)
 \quad\text{satisfies}\quad
 \mathcal R|_{E_v}\simeq\mathcal O_{E_v}
 \quad(v\in\mathbb Z^2).
\tag{5.17-A31}
\]
Indeed, the six toric curves detect \(\operatorname{Pic}(\mathrm{dP}_6)\):
in the blow-up basis \(H,E_1,E_2,E_3\), vanishing of the degrees on the
three exceptional curves kills the three \(E_i\)-coefficients, and
vanishing on any strict transform \(H-E_i-E_j\) kills the remaining
\(H\)-coefficient.  Since
\(H^1(\mathrm{dP}_6,\mathcal O)=0\), equality of those degrees implies
isomorphism, not merely numerical equivalence.

For \(b_\lambda=B_0\lambda\), the torus-translation part of
\(\Psi_\lambda\) preserves each toric prime divisor, while its fan part
translates the index by \(b_\lambda\).  Hence
\[
 \Psi_\lambda^*D_a-D_a
 =\sum_v(a(v+b_\lambda)-a(v))E_v
 =\operatorname{div}(\chi^{m_\lambda}),
\tag{5.17-A32}
\]
where
\[
 m_\lambda=(\ell_\lambda,k_\lambda),
 \qquad
 \ell_\lambda=B(b_\lambda,-),
 \qquad
 k_\lambda=a(b_\lambda).
\tag{5.17-A33}
\]
Thus multiplication by \(\chi^{m_\lambda}\) gives an individual
isomorphism
\(\Psi_\lambda^*\mathcal O(D_a)\simeq\mathcal O(D_a)\).
No cocycle assertion for these individual isomorphisms is needed: we use
them only to separate the horizontal winding of the genuine descent factor
of \(\mathcal M\).

For a ray \(\mathbb R_{\ge0}(v,1)\), its toric open set in
\(Y_\epsilon\) is
\[
 U_v\simeq\Delta_\epsilon\times(\mathbb C^*)^2,
 \qquad E_v^\circ=\{0\}\times(\mathbb C^*)^2.
\tag{5.17-A34}
\]
Indeed \((e_1,0),(e_2,0),(v,1)\) is an integral basis and \(t\) is the
character dual to the last vector; compare
\cite[\S2.1, pp.~28--29]{Fulton1993} and
\cite[Example~1.2.21, p.~31]{CoxLittleSchenck2011}.
The set \(U_v\) is Stein and retracts
onto \(E_v^\circ\).  By (5.17-A31),
\(c_1(\mathcal R|_{U_v})=0\); the exponential sequence and Cartan's
theorem~B therefore trivialize \(\mathcal R|_{U_v}\).  Choose a global
nowhere-zero frame \(s_v\) of \(\mathcal R|_{E_v}\).  A frame on \(U_v\) can be
multiplied by a unit pulled back from \((\mathbb C^*)^2\) so that its
restriction to \(E_v^\circ\) agrees with \(s_v|_{E_v^\circ}\); denote the
resulting frame by \(e_v\).

We also need to compare these frames on the generic torus.  The direction
of the toric affine-open inclusions is essential here.  If \(v,w\) are
adjacent vertices, choose a triangle \(\sigma=\{v,w,u\}\) containing them,
write \(z_v,z_w,z_u\) for the coordinates dual to its three height-one ray
generators, and put
\[
 S_\sigma:=U_\sigma\cap Y_\epsilon
 =\{(z_v,z_w,z_u)\in\mathbb C^3:
       |z_vz_wz_u|<\epsilon\}.
\tag{5.17-A34a}
\]
Because \(\rho_v<\sigma\), localization of the maximal affine chart gives
\begin{equation}
 U_{\rho_v}=D(z_wz_u)\subset U_\sigma,\qquad
 U_{\rho_w}=D(z_vz_u)\subset U_\sigma,\qquad
 U_{\rho_v}\cap U_{\rho_w}=D(z_vz_wz_u)=U_0.
\label{eq:analytic-toric-open-directions}
\end{equation}
Thus \(U_\sigma\) is not contained in \(U_{\rho_v}\cap U_{\rho_w}\), and
the coordinate discs in \(S_\sigma\) cannot be applied directly to the
ratio \(e_w/e_v\).

The space \(S_\sigma\) is nevertheless the required comparison chart.  It
is the closed graph
\[
 \{(z_v,z_w,z_u,t)\in\mathbb C^3\times\Delta_\epsilon:
   t=z_vz_wz_u\}
\]
in the Stein space \(\mathbb C^3\times\Delta_\epsilon\), and hence is
Stein.  Radial scaling of the three \(z\)-coordinates contracts it to the
origin.  The exponential sequence and Cartan's theorem~B therefore give
\(\operatorname{Pic}(S_\sigma)=0\).  Choose a frame \(e_\sigma\) of
\(\mathcal R|_{S_\sigma}\).

Let \(s_v\) be the global frame of
\(\mathcal R|_{E_v}\simeq\mathcal O_{E_v}\) used to normalize \(e_v\).
On \(U_{\rho_v}\cap Y_\epsilon\) put
\[
 A_{\sigma,v}=e_v/e_\sigma.
\]
At \(t=0\), this is the restriction to \((\mathbb C^*)^2\) of the
holomorphic unit
\[
 G_{\sigma,v}
 =\frac{s_v}{e_\sigma|_{E_v\cap S_\sigma}}
 \in\mathcal O^*(E_v\cap S_\sigma),
 \qquad E_v\cap S_\sigma=\{z_v=0\}\simeq\mathbb C^2.
\]
Consequently its winding on every horizontal cocharacter is zero at
\(t=0\).  Under
\(U_{\rho_v}\cap Y_\epsilon\simeq
\Delta_\epsilon\times(\mathbb C^*)^2\), that winding is an integer locally
constant in \(t\), and is therefore zero on every nearby nonzero fibre.
The same argument applies to \(A_{\sigma,w}\).  On their actual common
domain \(U_0\cap Y_\epsilon\),
\[
 e_w/e_v=A_{\sigma,w}/A_{\sigma,v},
\]
so this ratio has zero horizontal winding.  The one-skeleton of the
\(A_2\)-triangulation is connected, and multiplication along an edge path
now gives
\[
 \operatorname{wind}_m(e_w/e_v)=0
 \qquad(v,w,m\in\mathbb Z^2)
\tag{5.17-A34b}
\]
on every nonzero fibre.

Let
\(\theta^{\mathcal M}_\lambda:\Psi_\lambda^*\mathcal M\to\mathcal M\)
be the genuine descent isomorphism, let
\(\mu^+_\lambda:\Psi_\lambda^*\mathcal O(D_a)\to\mathcal O(D_a)\)
be multiplication by \(\chi^{m_\lambda}\) from (5.17-A32), and define the
inverse-dual isomorphism
\[
 \mu^-_\lambda:=((\mu^+_\lambda)^{-1})^\vee:
 \Psi_\lambda^*\mathcal O(-D_a)\longrightarrow\mathcal O(-D_a),
 \qquad s\longmapsto\chi^{-m_\lambda}s.
\]
The ordinary dual of \(\mu^+_\lambda\) has the opposite arrow; the
inverse-dual above is the map required by the pullback tensor product.  Then
\[
 \theta^{\mathcal R}_\lambda
 :=\theta^{\mathcal M}_\lambda\otimes\mu^-_\lambda:
 \Psi_\lambda^*\mathcal R\longrightarrow\mathcal R.
\]
Since \(\Psi_\lambda(U_v)=U_{v+b_\lambda}\), define
\(r_{\lambda,v}\in\mathcal O^*(U_v)\) by the typed identity
\begin{equation}
 \theta^{\mathcal R}_\lambda
   (\Psi_\lambda^*e_{v+b_\lambda})
 =r_{\lambda,v}e_v.
\label{eq:analytic-residual-descent-factor}
\end{equation}
At \(t=0\) it extends to an isomorphism between two trivial line bundles
on the compact surface \(E_v\).  Its factor in the chosen global frames is
therefore a nonzero constant.  The winding of
\(r_{\lambda,v}(t,-)\) around every horizontal one-parameter subgroup is
an integer locally constant in \(t\); at \(t=0\) it is zero, hence it is
zero for every nearby \(t\ne0\).

To express this isomorphism in one frame on the generic torus, put
\(q_{\lambda,v}=e_{v+b_\lambda}/e_v\).  Equation (5.17-A34b) gives zero
horizontal winding for \(q_{\lambda,v}\); because
\(\Psi_\lambda|_{Y_t}\) is multiplication by a fixed torus point, it
induces the identity on horizontal \(H_1\), and
\(\Psi_\lambda^*q_{\lambda,v}\) also has zero winding.  Relative to the
single frame \(e_v\), the residual multiplier is therefore
\[
 h_{\lambda,v}
 =\frac{r_{\lambda,v}}{\Psi_\lambda^*q_{\lambda,v}},
 \qquad
 \operatorname{wind}_m(h_{\lambda,v})=0.
\tag{5.17-A34c}
\]

On \(U_0\) let \(\mathbf1_a\) be the canonical toric frame of
\(\mathcal O(D_a)\), and put \(F_v=e_v\otimes\mathbf1_a\).  Equations
\cref{eq:analytic-residual-descent-factor} and (5.17-A32) give
\[
 \theta^{\mathcal M}_\lambda(\Psi_\lambda^*F_v)
 =h_{\lambda,v}\chi^{m_\lambda}F_v.
\]
On a nonzero fibre, \(\chi^{m_\lambda}=x^{\ell_\lambda}t^{k_\lambda}\).
With the section-descent convention (5.17am), the factor of automorphy is
therefore
\[
 \alpha_\lambda=(h_{\lambda,v}\chi^{m_\lambda})^{-1}.
\]
The residual factor has zero horizontal winding by (5.17-A34c), and hence
\begin{equation}
 \operatorname{wind}_m(\alpha_\lambda)
 =-\ell_\lambda(m)=-B(B_0\lambda,m).
\label{eq:analytic-factor-horizontal-winding}
\end{equation}

This winding is exactly the mixed Chern-class entry.  Pull \(F_v\) back to
the universal cover
\[
 p_t:\mathbb C^2\longrightarrow Y_t,
 \qquad \zeta\longmapsto(e^{2\pi i\zeta},t).
\]
The period indexed by \((\lambda,n)\in\mathbb Z^2\oplus\mathbb Z^2\)
acts by \(\zeta\mapsto\zeta+Z(s)\lambda+n\).  The horizontal factor is one,
because \(F_v\) is already defined on \(Y_t\).  Choose a logarithm
\(g_\lambda\in\mathcal O(\mathbb C^2)\) with
\[
 p_t^*\alpha_\lambda=e^{2\pi i g_\lambda}.
\]
For every horizontal period \(n\),
\[
 g_\lambda(\zeta+n)-g_\lambda(\zeta)
 =\operatorname{wind}_n(\alpha_\lambda).
\]
In the exponential connecting cocycle
\[
 c(l,l')=g_l(\zeta+l')+g_{l'}(\zeta)-g_{l+l'}(\zeta),
\]
take the horizontal logarithm to be zero and the logarithm for the combined
period \((\lambda,n)\) to be \(g_\lambda\).  Then
\[
 c((\lambda,0),(0,n))
 =\operatorname{wind}_n(\alpha_\lambda),
 \qquad
 c((0,n),(\lambda,0))=0.
\]
Therefore the alternating exponential connecting integer between
the vertical deck
translation indexed by \(\lambda\) and the horizontal loop indexed by
\(m\) is
\[
 E_M\bigl((\lambda,0),(0,m)\bigr)
 =-\ell_\lambda(m)
 =-B(B_0\lambda,m),
\]
which proves (5.17-A21).  This sign also recovers the explicit check in
\cref{prop:local-cusp-picard-obstruction}: for
\(a(v)=v^{\mathsf t}Hv\), one has \(B=2H\), so (5.17-A21) is the
cross-block \(-2B_0^{\mathsf t}H=-2P\), namely \(-2q\).

If \(E_M=0\), equation (5.17-A21) and the unimodularity of \(B_0\) force
\(B=0\).  The matrix in (5.17-A26) then gives
\(d_3=0\), \(d_2=0\), and \(d_1=0\).  Thus all six degrees of \(A\)
vanish, and the Picard detection argument following (5.17-A31) gives
\(A\simeq\mathcal O_{E_0}\).  This is (5.17-A22).
\end{proof}

\begin{theorem}[purely analytic vanishing of the global restriction]
\label{thm:analytic-all-picard}
\Verified\ For the constructed threefold,
\[
 r=\eta^*i^*:\operatorname{Pic}(X)\longrightarrow
 \operatorname{Pic}(\mathrm{dP}_6)
\]
is the zero homomorphism.  This proof uses only holomorphic
\(\mathcal O^*\)-factors, the two finite logarithmic fillings, the cusp
argument principle, and the \(\Psi\)-equivariant toric support function.
It uses neither a global integral Leray transgression nor local integral
nearby-cycle specialization.
\end{theorem}

\begin{proof}
Let \(L\in\operatorname{Pic}(X)\).  By
\cref{thm:coinvariants}, its invariant fibre class is
\[
 E=c_1(L|_F)=kq,
 \qquad k\in\mathbb Z,
 \qquad q=uw+6\gamma\delta.
\tag{5.17-A35}
\]
The exact contractions, now used without an integral specialization
lattice, are
\[
 \iota_{v_1}q=2\eta_1,
 \qquad
 \iota_{v_2}q=-3\eta_2,
 \qquad
 [\eta_1]=3[\delta],
 \qquad
 [\eta_2]=2[\delta].
\tag{5.17-A36}
\]
Apply (5.17-A5) with \((m_1,m_2)=(3,4)\).  In the real global
coinvariants,
\[
 \bar d_1=2k[\delta],
 \qquad
 \bar d_2=-\frac32k[\delta].
\tag{5.17-A37}
\]
Equations (5.17-A4) and (5.17-A6) therefore give
\[
 0=\bar d_1+\bar d_2+\bar d_0
   =\frac{k}{2}[\delta].
\tag{5.17-A38}
\]
Since \(V_G\otimes\mathbb R=\mathbb R[\delta]\), one has \(k=0\), so
\(E=0\).  Apply \cref{lem:analytic-support-specialization} to
\(M=L|_{N_0}\).  It gives
\[
 \eta^*i^*L=\nu^*(L|_W)\simeq\mathcal O_{E_0}.
\]
This holds for every \(L\), proving \(r=0\).
\end{proof}
